\documentclass[10pt]{article}
\usepackage[left=0.85in,right=0.85in,top=0.65in,bottom=0.65in]{geometry}
\usepackage{amsmath,amssymb,booktabs,array}
\usepackage[hidelinks]{hyperref}

\newcommand{\R}{\mathbb{R}}
\newcommand{\V}{\mathcal{V}}
\newcommand{\E}{\mathcal{E}}
\newcommand{\F}{\mathcal{F}}
\newcommand{\B}{\mathcal{B}}
\newcommand{\Sset}{\mathcal{S}}
\newcommand{\A}{\mathcal{A}}
\newcommand{\N}{\mathcal{N}}
\newcommand{\eps}{\varepsilon}

\title{Angle-Defect Reduction by Normal Displacement, Double-Cone Rays,
and Virtual-Sphere Rays}
\author{}
\date{}

\begin{document}
\maketitle
\vspace{-2em}

\section{Problem and notation}

Let a fixed-topology triangle mesh be
$\mathcal{M}=(\V,\E,\F)$ with vertex positions
$x_i\in\R^3$.  The integrated discrete Gaussian curvature is measured by
the angle defect
\begin{equation}
\Omega_i(x)=
\begin{cases}
2\pi-\displaystyle\sum_{f\ni i}\theta_{if}(x), & i\notin\B,\\[2mm]
\pi-\displaystyle\sum_{f\ni i}\theta_{if}(x), & i\in\B,
\end{cases}
\label{eq:defect}
\end{equation}
where $\B$ is the boundary-vertex set and $\theta_{if}$ is the corner angle
of face $f$ at vertex $i$.  For a tolerance $\eps>0$, an interior vertex is
classified as near-zero (yellow) when $|\Omega_i|\leq\eps$ and gray otherwise.

An initial singular set is fixed once,
\begin{equation}
\Sset=\{i\in\V\setminus\B:|\Omega_i(x^0)|>\tau\},
\qquad
\A=\V\setminus(\B\cup\Sset),
\end{equation}
with $\tau=1$ radian in the current experiments.  Boundary and singular
vertices are excluded from the threshold objective.  The principal curvature
objective is
\begin{equation}
E_{\mathrm{thr}}(x)=
\sum_{i\in\A}\left[\max(|\Omega_i(x)|-\eps,0)\right]^2.
\label{eq:threshold-energy}
\end{equation}
The topology is unchanged, so the total angle defect is checked after every
outer iteration as a Gauss--Bonnet consistency test.

\section{Local curvature evaluation}

Moving a set of vertices changes only triangles incident on that set and the
angle defects of vertices of those triangles.  If $\mathcal{H}$ is a monitored
vertex set, define
\begin{equation}
\F_{\mathcal H}=\{f\in\F:f\cap\mathcal H\neq\varnothing\}.
\end{equation}
Every face incident on a vertex in $\mathcal H$ belongs to
$\F_{\mathcal H}$; consequently, recomputing corner angles only on
$\F_{\mathcal H}$ gives the exact new defects for all vertices in
$\mathcal H$.  This replaces repeated full-mesh evaluations inside an
optimizer.  Full-mesh defects and geometry checks are still recomputed before
accepting an outer update and at the end of each iteration.

For a single-vertex ray candidate, only faces incident on vertex $i$ change.
If $\Delta\theta_{jf}=\theta_{jf}(x)-\theta_{jf}(x')$, then
\begin{equation}
\Omega_j(x')=\Omega_j(x)+
\sum_{\substack{f\ni i\\f\ni j}}\Delta\theta_{jf}.
\label{eq:local-defect-update}
\end{equation}
Equation~\eqref{eq:local-defect-update} is used for the cone and sphere
candidate evaluations.  Its numerical result agrees with a full recomputation
to floating-point precision.

\section{Method I: local-objective normal displacement}

At outer iteration $k$, define the moderate-gray target set
\begin{equation}
\mathcal{G}^k=\{i\in\A:\eps<|\Omega_i(x^k)|\leq\tau\}.
\end{equation}
The movable set contains the targets and $r$ vertex rings,
\begin{equation}
\mathcal{M}^k=\N_r(\mathcal{G}^k)\setminus(\B\cup\Sset),
\qquad
\mathcal{H}^k=\N_1(\mathcal{M}^k).
\end{equation}
The current implementation uses $r=1$.  Vertex normals $n_i^k$ are frozen
during one inner optimization, and each movable position is parameterized by
one scalar:
\begin{equation}
x_i(\delta)=x_i^k+\delta_i n_i^k,
\qquad i\in\mathcal{M}^k.
\label{eq:normal-param}
\end{equation}
Let $\ell_i^k$ be the current average incident-edge length and $\ell_i^0$ its
initial value.  The box bound combines a per-step and cumulative displacement
budget:
\begin{equation}
|\delta_i|\leq
\min\left(
\alpha\ell_i^k,
\beta\ell_i^0-\|x_i^k-x_i^0\|
\right),
\label{eq:normal-bound}
\end{equation}
where negative remaining budgets are clamped to zero.

The inner objective is
\begin{equation}
E_{\mathrm{normal}}=
E_\kappa+\lambda_0E_0+\lambda_sE_s+
\lambda_pE_p+\lambda_lE_l+\lambda_eE_e.
\label{eq:normal-objective}
\end{equation}
For monitored active vertices $\mathcal{H}_A=\mathcal{H}^k\cap\A$,
\begin{align}
E_\kappa &=\frac{1}{|\mathcal{H}_A|}
\sum_{i\in\mathcal{H}_A}
\left(\frac{\max(|\Omega_i|-\eps,0)}{\eps}\right)^2,\\
E_0 &=\frac{1}{|\mathcal{H}_A|}
\sum_{i\in\mathcal{H}_A}\left(\frac{\Omega_i}{\eps}\right)^2,\\
E_s &=\frac{1}{|\mathcal{H}^k\cap\Sset|}
\sum_{i\in\mathcal{H}^k\cap\Sset}
\left(\frac{\Omega_i-\Omega_i^0}{\max(|\Omega_i^0|,\eps)}\right)^2,\\
E_p &=\frac{1}{|\mathcal{M}^k|}
\sum_{i\in\mathcal{M}^k}\left(\frac{\delta_i}{\ell_i^k}\right)^2.
\end{align}
Let $d_i=\delta_i n_i^k$ for movable vertices and $d_i=0$ otherwise.
For the edges affected by $\mathcal{M}^k$,
\begin{align}
E_l &= \frac{1}{|\E_M|}\sum_{(i,j)\in\E_M}
\frac{\|d_i-d_j\|^2}{[\tfrac12(\ell_i^k+\ell_j^k)]^2},\\
E_e &= \frac{1}{|\E_M|}\sum_{(i,j)\in\E_M}
\left(\frac{\|x_i(\delta)-x_j(\delta)\|-\|x_i^k-x_j^k\|}
{\|x_i^k-x_j^k\|}\right)^2.
\end{align}
The current default weights are
$\lambda_0=10^{-3}$, $\lambda_s=0.1$, $\lambda_p=0.01$,
$\lambda_l=0.05$, and $\lambda_e=0.1$.

The box-constrained problem is solved by L-BFGS-B.  The Bunny experiment uses
five inner iterations per outer iteration.  Backtracking scales the proposed
update by $2^{-q}$, $q=0,\ldots,10$.  An update is accepted only if the
full-mesh $E_{\mathrm{thr}}$ decreases, all triangles pass the per-update
orientation and area tests, affected edges pass the per-update distortion
limit, cumulative vertex displacement satisfies~\eqref{eq:normal-bound}, the
original-relative edge distortion is below $0.25$, the original-relative area
ratio is above $0.30$, and all defects remain finite.

\section{Method II: double-cone ray search}

The ray methods update moderate-gray vertices sequentially in decreasing order
of $|\Omega_i|$.  At vertex $i$, construct an orthonormal frame
$(n_i,t_{i1},t_{i2})$.  For an outward-cone sample with azimuth $\phi_j$ and
$c_j\in[\cos\vartheta,1]$, define
\begin{equation}
r_j^+=c_jn_i+\sqrt{1-c_j^2}
(\cos\phi_j\,t_{i1}+\sin\phi_j\,t_{i2}),
\qquad r_j^-=-r_j^+.
\label{eq:cone-rays}
\end{equation}
The implementation samples $c_j$ uniformly over the spherical cap and uses a
golden-angle sequence for $\phi_j$.  The retained experiment uses 30 total rays
and cone half-angle $\vartheta=60^\circ$, corresponding to a full opening angle
of $120^\circ$ for each cone.

The search interval on every ray is
\begin{equation}
0\leq t\leq L_i,
\qquad
L_i=\min\left(\gamma\ell_i^k,
\beta\ell_i^0-\|x_i^k-x_i^0\|\right),
\label{eq:ray-bound}
\end{equation}
and a candidate is $x_i(t)=x_i^k+tr_j$.  The current parameters are
$\gamma=0.1$ and $\beta=0.1$.  Five equally spaced coarse intervals are
tested.  If two consecutive samples bracket a sign change of $\Omega_i(t)$,
binary search refines the bracket for at most 12 steps or until
$|\Omega_i|\leq10^{-5}$.

Candidates are ranked by
\begin{equation}
E_{\mathrm{ray}}(x)=E_{\mathrm{thr}}(x)+\lambda_sE_s(x),
\label{eq:ray-merit}
\end{equation}
using the exact local update in~\eqref{eq:local-defect-update}.  A candidate is
accepted only if both $E_{\mathrm{ray}}$ and $E_{\mathrm{thr}}$ decrease and
the incident triangles, incident edges, and cumulative vertex displacement
pass the geometry-safety checks.  The same original-relative edge-distortion
and area-ratio bounds used by the normal method are enforced.  A full defect
computation is performed after each outer iteration.

The cone construction is used only as a direction sampler.  The method does
not compute the Shape Diameter Function and does not use a thickness value.

\section{Method III: virtual-sphere ray search}

The virtual-sphere method uses the same search interval, coarse sampling,
binary search, merit function, sequential update, and safety checks as the
double-cone method.  Only the direction set changes.  For $m$ rays and
$j=0,\ldots,m-1$, Fibonacci-sphere directions are
\begin{align}
z_j &= 1-\frac{2(j+1/2)}{m},\\
\phi_j &= j\pi(3-\sqrt5),\\
r_j &= \left(\sqrt{1-z_j^2}\cos\phi_j,
\sqrt{1-z_j^2}\sin\phi_j,z_j\right).
\label{eq:sphere-rays}
\end{align}
Thus the virtual sphere samples unrestricted directions, including normal,
tangential, and oblique motion.  The current experiment uses $m=30$.

\section{Stanford Bunny results at \texorpdfstring{$\eps=0.1$}{epsilon=0.1}}

The full-resolution mesh has 34,834 vertices, 69,451 triangles, 223 boundary
vertices, and seven fixed initial singular vertices.  Boundary vertices are
excluded from the reported interior-gray counts and threshold energy.

\begin{center}
\begin{tabular}{lrrrr}
\toprule
Method & Moderate gray & $E_{\mathrm{thr}}$ & Yellow BFS region & Max. edge change\\
\midrule
Original & 533 & 14.1534 & 34,136 & 0\%\\
Local normal & \textbf{156} & \textbf{0.6757} & \textbf{34,518} & 24.71\%\\
Double cone & 202 & 3.6586 & 34,471 & 19.14\%\\
Virtual sphere & 213 & 3.7824 & 34,461 & 23.97\%\\
\bottomrule
\end{tabular}
\end{center}

The local normal method gives the strongest curvature reduction, but it also
approaches the accumulated edge-distortion bound and reduces the minimum triangle
quality from $0.01390$ to $0.00841$.  Its per-update orientation test does not
guarantee that every final face normal remains aligned with the original face
normal.  The cone and sphere methods preserve the original minimum triangle
quality in this experiment; the cone produces fewer gray vertices, whereas the
sphere retains a larger minimum triangle-area ratio ($0.594$ versus $0.313$).

\section{Current limitations and required safeguards}

All three methods now enforce an original-relative accumulated edge bound
\begin{equation}
\max_{(i,j)\in\E}
\frac{\big|\|x_i-x_j\|-\|x_i^0-x_j^0\|\big|}{\|x_i^0-x_j^0\|}
\leq\rho_{\mathrm{cum}},
\end{equation}
with $\rho_{\mathrm{cum}}=0.25$, together with an original-relative area-ratio
bound $A_f/A_f^0\geq0.30$.  These guards control extreme accumulated changes.
The final geometry comparison also reports RMS displacement, 95th-percentile
affected-edge distortion, 5th-percentile affected-face area ratio, triangle
quality, and face-normal deviation.  A future version may add triangle-quality
and face-normal limits.  A curvature-transfer term should also penalize creating
new gray vertices while correcting existing targets.

\end{document}
